% Part: proof-theory
% Chapter: natural-deduction
% Section: rules-proofs

\documentclass[../../../include/open-logic-section]{subfiles}

\begin{document}

\olfileid{pt}{nat}{rp}

\olsection{في القواعد و\usetoken{P}{proof}}

نقدم أولًا تعريفات نثبت بها الاصطلاح.

ولننظر في قاعدتي~\Log{N1c} و\Log{N1i} اللتين يرد فيهما
$!A \lif !B$:
\[
\bottomAlignProof
\AxiomC{$\Discharge{!A}{{\OLId{latin-ordinary}{x}}}$}
\shortDeduce
\DeduceC{$!B$}
\DischargeRule{\Intro{\lif}}{{\OLId{latin-ordinary}{x}}}
\UnaryInfC{$!A \lif !B$}
\DisplayProof
\qquad
\bottomAlignProof
\AxiomC{$!A \lif !B$}
\AxiomC{$!A$}
\RightLabel{\Elim{\lif}}
\BinaryInfC{$!B$}
\DisplayProof
\]
أما قاعدة~$\Elim{\lif}$ فأمرها بيّن. فالـ!!{formula}s التي فوق الخط
تسمى \emph{المقدمات}، والـ!!{formula} الواقعة عن اليسار، وهي
!!{formula}~$!A \lif !B$، تسمى \emph{الرئيسة}. ولكل قاعدة
!!{elimination} مقدمة من هذا القبيل، موضعها دائمًا أقصى اليسار، وهي
\emph{المقدمة الكبرى} للاستدلال. وما زاد عليها، كما في مثالنا، فهو
\emph{مقدمات صغرى}. ومقدمات قواعد~!!{introduction} تسمى أيضًا مقدمات
كبرى. وأما~!!{formula} الواقعة تحت الخط فهي \emph{النتيجة}.

ولقاعدة~$\Intro{\lif}$ مقدمة واحدة، ونتيجتها هنا هي الصيغة الرئيسة
$!A \lif !B$. وعلى هذا تجري قواعد~!!{introduction} كلها: فالنتيجة
هي الصيغة الرئيسة، وعاملها الرئيس هو العامل «المدخل».

وبيان الرمز~$\Discharge{!A}{{\OLId{latin-ordinary}{x}}}$ كما يأتي. يطبق، في~!!a{proof}، استدلال
$\Intro{\lif}$ على~!!a{formula}~$!B$ مقدمةً له. وللـ!!{proof} الجزئي
المنتهي بـ$!B$~!!{formula}s أول، نسميها \emph{فروضًا}. فإذا كان
!!^a{proof} لـ$!B$ قائمًا على فرض أو أكثر من الصورة~$!A$، كان~!!a{proof}
على أن~$!A$ تستلزم~$!B$. ثم إذا أعملنا قاعدة~$\Intro{\lif}$ استنتجنا
$!A \lif !B$، وانقطع اعتماد النتيجة على الفرض~$!A$. فنقول، دلالة على
ذلك، إن الفروض من الصورة~$!A$ قد \emph{برئت} أو \emph{أغلقت} في
!!{proof} المشتمل على استدلال~$\Intro{\lif}$. والقاعدة هي التي تعين ما
يجوز تبرئته؛ ففي~$\Intro{\lif}$ التي نتيجتها~$!A \lif !B$ تكون الفروض
على الصورة~$!A$، أي توافق مقدم الشرط. وقد نحتاج إلى تعيين الفروض التي
تبرئها كل قاعدة، ولذلك وضعت علامة التبرئة~${\OLId{latin-ordinary}{x}}$. غير أن القاعدة المبرئة
لا \emph{يلزمها} أن تبرئ \emph{كل} فرض من الصورة المطلوبة، بل يجوز ألا
تبرئ فرضًا. فالمثال الآتي صحيح، مع أن~$\Intro{\lif}$ الأولى لا تبرئ
شيئًا:
\[
\AxiomC{$\Discharge{!A}{{\OLId{latin-ordinary}{y}}}$}
\DischargeRule{\Intro{\lif}}{{\OLId{latin-ordinary}{x}}}
\UnaryInfC{$!B \lif !A$}
\DischargeRule{\Intro{\lif}}{{\OLId{latin-ordinary}{y}}}
\UnaryInfC{$!A \lif (!B \lif !A)$}
\DisplayProof
\]
وإذا لم تبرئ القاعدة فرضًا أسقطنا علامة التبرئة، كما في هذا الموضع.

في~!!{proof}
\[
\AxiomC{$\Discharge{!A}{{\OLId{latin-ordinary}{x}}}$}
\AxiomC{$\Discharge{!A}{{\OLId{latin-ordinary}{y}}}$}
\RightLabel{\Intro{\land}}
\BinaryInfC{$!A \land !A$}
\RightLabel{\Elim{\land}}
\UnaryInfC{$!A$}
\DischargeRule{\Intro{\lif}}{{\OLId{latin-ordinary}{x}}}
\UnaryInfC{$!A \lif !A$}
\DischargeRule{\Intro{\lif}}{{\OLId{latin-ordinary}{y}}}
\UnaryInfC{$!A \lif (!A \lif !A)$}
\DisplayProof
\]
يبرئ استدلال~$\Intro{\lif}$ الأول الفرض الأيسر~$!A^{\OLId{latin-ordinary}{x}}$، ويبرئ الثاني
الفرض الأيمن~$!A^{\OLId{latin-ordinary}{y}}$. أي إن كل الفروض الموسومة~${\OLId{latin-ordinary}{x}}$ يبرئها
$\Intro{\lif}$ ذو العلامة~${\OLId{latin-ordinary}{x}}$، وكل الفروض الموسومة~${\OLId{latin-ordinary}{y}}$ يبرئها
الاستدلال الموسوم~${\OLId{latin-ordinary}{y}}$. ولكي ينجح ذلك، من المهم أن تكون كل الفروض ذات
العلامة الواحدة~!!{formula} نفسها كذلك.

وأما قواعد الأسوار ففيها زيادة تعقيد. ومن ذلك قواعد~$\lexists$ في
\Log{N1i}:
\[
\bottomAlignProof
\AxiomC{$!A({\OLId{latin-ordinary}{t}})$}
\RightLabel{\Intro{\lexists}}
\UnaryInfC{$\lexists[{\OLId{latin-ordinary}{x}}][!A({\OLId{latin-ordinary}{x}})]$}
\DisplayProof
\qquad
\bottomAlignProof
\AxiomC{$\lexists[{\OLId{latin-ordinary}{x}}][!A({\OLId{latin-ordinary}{x}})]$}
\AxiomC{$\Discharge{!A({\OLId{latin-ordinary}{c}})}{{\OLId{latin-ordinary}{x}}}$}
\shortDeduce
\DeduceC{$!B$}
\DischargeRule{\Elim{\lexists}}{{\OLId{latin-ordinary}{x}}}
\BinaryInfC{$!B$} \DisplayProof
\bottomAlignProof
\]
مقدمة~$\Intro{\lexists}$ هي~$!A({\OLId{latin-ordinary}{t}})$، على أن يكون~${\OLId{latin-ordinary}{t}}$ حدًا مغلقًا، أي
لا متغير فيه. ويصح الاستدلال، أي يطابق مخطط~$\Intro{\lexists}$، إذا
وجدت~!!{formula}~$!A$ ومتغير~${\OLId{latin-ordinary}{x}}$ وحد مغلق~${\OLId{latin-ordinary}{t}}$، وكانت النتيجة
$\lexists[{\OLId{latin-ordinary}{x}}][!A]$ والمقدمة~$\Subst{!A}{{\OLId{latin-ordinary}{t}}}{{\OLId{latin-ordinary}{x}}}$. وأما
$\Elim{\lexists}$ فتبرئ فروضًا من الصورة~$!A({\OLId{latin-ordinary}{c}})$؛ فلكل استدلال منها
!!{constant}~${\OLId{latin-ordinary}{c}}$، والفروض الجائز تبرئتها من الصورة
$\Subst{!A}{{\OLId{latin-ordinary}{c}}}{{\OLId{latin-ordinary}{x}}}$. ويجب أن تتفق هذه الفروض كلها في الثابت~${\OLId{latin-ordinary}{c}}$. وهذا
الثابت هو \emph{المتغير المميز} لاستدلال~$\Elim{\lexists}$. وإنما يصح
الاستدلال إذا لم يظهر~${\OLId{latin-ordinary}{c}}$ في فرض مفتوح غير~$!A({\OLId{latin-ordinary}{c}})$ المبرأ، ولا في~$!B$،
ولا في~$!A$؛ وهذا هو \emph{شرط المتغير المميز}.

!!^{proof}s في~\Log{N1c} و\Log{N1i} أشجار منتهية من وقائع الصيغ، تُولّد
استقرائيًا من الفروض بتطبيق قواعد الاستدلال. وكل ورقة فرض، قد تكون
موسومة بعلامة تبرئة، وكل واقعة~!!{formula} أخرى تلزم من~!!{formula}s
الواقعة فوقها مباشرة بإحدى قواعد استدلال النظام. وإذا كانت واقعةُ صيغةٍ
فرضًا في الجزء من الشجرة الواقع فوق استدلال، ولم تكن قد بُرئت باستدلال
يقع فوق ذلك الاستدلال، سميت فرضًا مفتوحًا عند ذلك الاستدلال.

\begin{defn}\ollabel{defn:proof}
تُعرّف \emph{!!^{proof}s} في~\Log{N1c} و\Log{N1i} استقرائيًا بوصفها
أشجارًا منتهية من~!!{formula}s كما يأتي:
\begin{enumerate}
  \item\ollabel{defn:prf:assumption} كل~!!{formula} موسومة، بمفردها،
  بعلامة تبرئة مثل~${\OLId{latin-ordinary}{x}}$ أو~${\OLId{latin-ordinary}{y}}$ هي~!!a{proof}. وفي~!!a{proof} يتألف من
  !!{formula} واحدة، تُعد~!!{formula} فرضًا مفتوحًا.
  \item\ollabel{defn:prf:one-prem} إذا كانت~${\OLId{latin-looped}{R}}$ قاعدة ذات مقدمة أو
  مقدمتين أو ثلاث مقدمات~$!A_{\OLMathNumeral{1}}$ و$!A_{\OLMathNumeral{2}}$ و$!A_{\OLMathNumeral{3}}$، ونتيجتها~$!A$، وكانت
  $\delta_{\OLMathNumeral{1}}$ و$\delta_{\OLMathNumeral{2}}$ و$\delta_{\OLMathNumeral{3}}$~!!{proof}s تنتهي بـ$!A_{\OLMathNumeral{1}}$ و
  $!A_{\OLMathNumeral{2}}$ و$!A_{\OLMathNumeral{3}}$ على الترتيب، فإن كلًا مما يأتي
  \[
  \AxiomC{}
  \RightLabel{$\delta_{\OLMathNumeral{1}}$}
  \DeduceC{$!A_{\OLMathNumeral{1}}$}
  \RightLabel{${\OLId{latin-looped}{R}}$}
  \UnaryInfC{$!A$}
  \DisplayProof
\qquad
  \AxiomC{}
  \RightLabel{$\delta_{\OLMathNumeral{1}}$}
  \DeduceC{$!A_{\OLMathNumeral{1}}$}
  \AxiomC{}
  \RightLabel{$\delta_{\OLMathNumeral{2}}$}
  \DeduceC{$!A_{\OLMathNumeral{2}}$}
  \RightLabel{${\OLId{latin-looped}{R}}$}
  \BinaryInfC{$!A$}
  \DisplayProof
\qquad
  \AxiomC{}
  \RightLabel{$\delta_{\OLMathNumeral{1}}$}
  \DeduceC{$!A_{\OLMathNumeral{1}}$}
  \AxiomC{}
  \RightLabel{$\delta_{\OLMathNumeral{2}}$}
  \DeduceC{$!A_{\OLMathNumeral{2}}$}
  \AxiomC{}
  \RightLabel{$\delta_{\OLMathNumeral{3}}$}
  \DeduceC{$!A_{\OLMathNumeral{3}}$}
  \RightLabel{${\OLId{latin-looped}{R}}$}
  \TrinaryInfC{$!A$}
  \DisplayProof
  \]
  هو~!!a{proof}، على الترتيب، بشرط توافق علامات الفروض في~!!{proof}s،
  أي ألا تحمل~!!{formula}s مختلفة علامة التبرئة نفسها، وأن تكون
  الاستدلالات تطبيقات صحيحة لـ${\OLId{latin-looped}{R}}$.
  \item تُعد أعلى وقائع~!!{formula} في~!!{proof} الجديد
  \emph{فروضًا مفتوحة} لـ!!{proof} إذا كانت فروضًا مفتوحة في
  $\delta_{\OLMathNumeral{1}}$ أو$\delta_{\OLMathNumeral{2}}$ أو$\delta_{\OLMathNumeral{3}}$، باستثناء الوقائع التي تبرئها
  القاعدة. فإذا وُسمت القاعدة بـ${\OLId{latin-ordinary}{x}}$، أو بـ${\OLId{latin-ordinary}{x}}\,{\OLId{latin-ordinary}{y}}$، كانت الوقائع المبرأة:
  في~$\Intro{\lif}$ و$\Intro{\lnot}$ كل الفروض المفتوحة في~$\delta_{\OLMathNumeral{1}}$
  الموسومة~${\OLId{latin-ordinary}{x}}$؛ وفي~$\Elim{\lexists}$ كل الفروض المفتوحة الموسومة~${\OLId{latin-ordinary}{x}}$
  في~$\delta_{\OLMathNumeral{2}}$؛ وفي~$\Elim{\lor}$ كل الفروض المفتوحة الموسومة~${\OLId{latin-ordinary}{x}}$ في
  $\delta_{\OLMathNumeral{2}}$ وكل الفروض المفتوحة الموسومة~${\OLId{latin-ordinary}{y}}$ في~$\delta_{\OLMathNumeral{3}}$.
\end{enumerate}
\end{defn}

تسمى~!!{proof}s المستخدمة في البناء الاستقرائي لـ!!a{proof}
\emph{!!{proof}s الجزئية} له. وتسمى~!!{proof}s الجزئية المنتهية بمقدمات
استدلال~\emph{!!{proof}s الجزئية المباشرة} للاستدلال. وترد قواعد
\Log{N1c} و\Log{N1i} في~\olref{tab:N1}.

\subfile{rules-N1}

\begin{defn}
يُعرّف \emph{!!{height}}~$\pheight{\delta}$ لـ!!a{proof}~$\delta$
استقرائيًا كما يأتي.
\begin{enumerate}
  \item إذا تألف~$\delta$ من فرض، كان~$\pheight{\delta}={\OLMathNumeral{0}}$.
  \item إذا انتهى~$\delta$ باستدلال ذي مقدمة واحدة و!!{proof} جزئي
  مباشر~$\delta_{\OLMathNumeral{1}}$، كان~$\pheight{\delta}=\pheight{\delta_{\OLMathNumeral{1}}}+{\OLMathNumeral{1}}$.
  \item إذا انتهى~$\delta$ باستدلال ذي مقدمتين و!!{proof}s جزئية
  مباشرة~$\delta_{\OLMathNumeral{1}}$ و$\delta_{\OLMathNumeral{2}}$، كان~$\pheight{\delta}=
  \max(\pheight{\delta_{\OLMathNumeral{1}}},\pheight{\delta_{\OLMathNumeral{2}}})+{\OLMathNumeral{1}}$.
  \item إذا انتهى~$\delta$ بـ$\Elim{\lor}$ و!!{proof}s جزئية مباشرة
  $\delta_{\OLMathNumeral{1}}$ و$\delta_{\OLMathNumeral{2}}$ و$\delta_{\OLMathNumeral{3}}$، كان~$\pheight{\delta}=
  \max(\pheight{\delta_{\OLMathNumeral{1}}},\pheight{\delta_{\OLMathNumeral{2}}},\pheight{\delta_{\OLMathNumeral{3}}})+{\OLMathNumeral{1}}$.
\end{enumerate}
إذا كان للمتتالية~${\OLId{latin-looped}{S}}$~!!a{proof} ارتفاعه~$\le {\OLId{latin-ordinary}{n}}$، كتبنا~$\Proves[{\OLId{latin-ordinary}{n}}] {\OLId{latin-looped}{S}}$.
\end{defn}

\begin{ex}
في~\Log{N1i}، تُعد أي~!!{formula}~!!a{proof} من نفسها بوصفها فرضًا
مفتوحًا. ومن ثم
\[
!A \land !B^{\OLId{latin-ordinary}{x}}
\]
هو~!!a{proof}~$\delta_{\OLMathNumeral{1}}$، وارتفاعه~${\OLMathNumeral{0}}$.
ومن~$\delta_{\OLMathNumeral{1}}$ نستطيع الحصول على~!!{proof}s اثنين~$\delta_{\OLMathNumeral{2}}$ و
$\delta_{\OLMathNumeral{3}}$ بتطبيق إحدى نسختي~$\Elim{\land}$:
\[
\AxiomC{$!A \land !B^{\OLId{latin-ordinary}{x}}$}
\RightLabel{\Elim{\land}[{\OLMathNumeral{2}}]}
\UnaryInfC{$!B$}
\DisplayProof
\qquad
\AxiomC{$!A \land !B^{\OLId{latin-ordinary}{x}}$}
\RightLabel{\Elim{\land}[{\OLMathNumeral{1}}]}
\UnaryInfC{$!A$}
\DisplayProof
\]
والـ!!{proof} الجزئي المباشر للاستدلال الأخير في كل من~$\delta_{\OLMathNumeral{2}}$ و
$\delta_{\OLMathNumeral{3}}$ هو~$\delta_{\OLMathNumeral{1}}$. وفي كل~!!{proof} منهما تكون واقعة
$!A\land !B$ فرضًا مفتوحًا. ولدينا
$\pheight{\delta_{\OLMathNumeral{2}}}=\pheight{\delta_{\OLMathNumeral{3}}}={\OLMathNumeral{1}}$.

تتطلب قاعدة~$\Intro{\land}$ مقدمتين من الصورتين~$!B$ و$!A$ لتسويغ
نتيجة من الصورة~$!B \land !A$. ولذلك يكون الآتي~!!a{proof}~$\delta_{\OLMathNumeral{4}}$:
\[
\AxiomC{$!A \land !B^{\OLId{latin-ordinary}{x}}$}
\RightLabel{\Elim{\land}[{\OLMathNumeral{2}}]}
\UnaryInfC{$!B$}
\LeftSubproofLabel{$\delta_{\OLMathNumeral{2}}$}
\AxiomC{$!A \land !B^{\OLId{latin-ordinary}{x}}$}
\RightLabel{\Elim{\land}[{\OLMathNumeral{1}}]}
\UnaryInfC{$!A$}
\RightSubproofLabel{$\delta_{\OLMathNumeral{3}}$}
\RightLabel{\Intro{\land}}
\BinaryInfC{$!B \land !A$}
\DisplayProof
\]
وارتفاعه~$\pheight{\delta_{\OLMathNumeral{4}}}=\max(\pheight{\delta_{\OLMathNumeral{2}}},
\pheight{\delta_{\OLMathNumeral{3}}})+{\OLMathNumeral{1}}=\max({\OLMathNumeral{1}},{\OLMathNumeral{1}})+{\OLMathNumeral{1}}={\OLMathNumeral{2}}$. وله واقعتان من
$!A \land !B$ بوصفهما فرضين، وكلتاهما مفتوحة.

وأخيرًا نستطيع تطبيق~$\Intro{\lif}$ للحصول على~$\delta_{\OLMathNumeral{5}}$:
\[
\AxiomC{$\Discharge{!A \land !B}{{\OLId{latin-ordinary}{x}}}$}
\RightLabel{\Elim{\land}[{\OLMathNumeral{2}}]}
\UnaryInfC{$!B$}
\LeftSubproofLabel{$\delta_{\OLMathNumeral{2}}$}
\AxiomC{$\Discharge{!A \land !B}{{\OLId{latin-ordinary}{x}}}$}
\RightLabel{\Elim{\land}[{\OLMathNumeral{1}}]}
\UnaryInfC{$!A$}
\RightSubproofLabel{$\delta_{\OLMathNumeral{3}}$}
\RightLabel{\Intro{\land}}
\BinaryInfC{$!B \land !A$}
\DischargeRule{\Intro{\lif}}{{\OLId{latin-ordinary}{x}}}
\UnaryInfC{$(!A \land !B) \lif (!B \land !A)$}
\DisplayProof
\]
وارتفاعه~$\pheight{\delta_{\OLMathNumeral{4}}}+{\OLMathNumeral{1}}={\OLMathNumeral{3}}$. ويبرئ استدلال~$\Intro{\lif}$ كل
وقائع الفروض من الصورة~$!A \land !B$ الموسومة~${\OLId{latin-ordinary}{x}}$، أي الفرضين اللذين
كانا مفتوحين من قبل في البرهان الجزئي المباشر. ولأنهما الفرضان
الوحيدان لـ$\delta_{\OLMathNumeral{5}}$، فلا تكون له فروض مفتوحة، ومن ثم يكون~!!a{proof}
\emph{لـ}!!{formula} نهايته~$(!A \land !B) \lif (!B \land !A)$.
\end{ex}

\end{document}
